\section{Pilot Selective Proof Rail}
\label{sec:pilot-mobile-selective-proof}

The unified Pilot product reuses GlyphChain as an optional tamper-evident proof
rail, not as a store for private life or business data and not as a payment rail.
The authoritative message, task, file, contact, calendar and memory records remain
inside the customer-controlled mother brain.  A proof commitment can demonstrate
that a bounded event occurred without disclosing the event content.

\subsection{Versioned commitment model}

The protocol \texttt{pilot.selective-proof-rail.v1} defines five event classes:
authorization, delivery, acceptance, outcome and deletion.  A public commitment is
limited by construction to a version, commitment identifier, opaque object
reference, canonical scope hash, policy version, timestamp, bounded outcome,
source-receipt hash and optional links to an earlier commitment.  The mother signs
this exact projection with its Ed25519 device identity.

The opaque object reference is derived from the local identifier rather than
publishing that identifier.  Message text, voice recordings, files, names, contact
routes, private memory, calendar details, location, credentials, account tokens and
payment information are not accepted fields in the public structure.  The supplied
mother-private record is hashed canonically and is never copied into the proof-rail
state.

\subsection{Task lifecycle integration}

For a structured Pilot-to-Pilot task, the sender mother emits an authorization
commitment after exact phone approval and a delivery commitment after constructing
the encrypted recipient-mother packet.  The recipient mother emits a separate
acceptance or decline commitment from its signed response receipt.  Completion,
failure, expiry and other terminal task outcomes can create outcome commitments.

These commitments do not redefine application state.  Delivery means delivery to
the recipient mother, not human reading.  Acceptance requires the recipient's
signed action.  No proof may convert a pending action into a sent, accepted, paid or
completed claim.  The encrypted task packet and its ordinary signed receipts remain
the operational authority.

\subsection{Idempotency, correction and deletion}

Every emission requires a bounded idempotency key.  Repeating the same operation
returns the original commitment; reusing the key for different content is rejected.
Corrections append a new commitment linked through \texttt{correction\_of} and
\texttt{previous\_commitment\_id}; history is not rewritten.

When an authorized private source record is deleted, Pilot can append a deletion
tombstone linked to the earlier proof.  The non-identifying historical hash may
remain immutable, while the private record and any ability to reproduce its content
are removed under the customer's retention policy.  Pilot must explain this
distinction before deletion rather than implying that an already published hash can
be erased.

\subsection{Verification and availability policy}

An authorized mother verifies a local source record by recomputing its canonical
hash and checking the mother signature over the public commitment.  Any change to
the private receipt or signed commitment fails verification.  Verification without
the private source can establish commitment integrity but cannot reconstruct or
disclose the source content.

Proof publication has three explicit policies.  \emph{Local only} keeps the signed
commitment on the mother.  \emph{Best effort} continues the permitted application
operation and marks publication pending when GlyphChain is unavailable.
\emph{Required} fails closed before claiming publication.  Application health,
chain health and pending-publication count are reported separately, preventing a
chain outage from being confused with message or task delivery.

\subsection{Reuse of the existing GlyphChain foundation}

The adapter maps the allowlisted Pilot commitment to the existing AION evidence
proof envelope and commit/lookup store.  Only the privacy-safe projection enters
that envelope.  The existing canonical codec, commitment verification and receipt
mechanism are reused; wallet, PHO, token, payment and escrow features are neither
required nor enabled.  Every returned safety record explicitly reports those
economic functions as false.

\subsection{Verification evidence}

Automated tests cover allowlist inspection, absence of private strings, duplicate
emission, malicious idempotency-key reuse, modified private receipts, modified
signed commitments, append-only correction, deletion tombstones, offline
best-effort behavior, required-publication failure and the existing GlyphChain
adapter.  An end-to-end two-mother task test produces authorization, delivery and
acceptance commitments while confirming that the task title is absent from both
proof stores.  No wallet, balance, token transfer, payment or escrow is created.
